\chapter{Einführung}
\label{chap:einfuehrung}

Dies ist das erste Kapitel des Buches. Die Dateien sind so organisiert, dass jedes
Kapitel einzeln unter \texttt{kapitel/} bearbeitet werden kann. Ein Push auf den
Branch \texttt{main} stößt den automatischen PDF-Build an.

\section{Schreiben mit LaTeX}

LaTeX übernimmt die Formatierung von Überschriften, Tabellen und Literaturangaben.
Eine Einführung in das Satzsystem bietet \textcite{lamport1994}.
Innerhalb des Dokuments lassen sich Abschnitte wie
\autoref{sec:beispiel} referenzieren.

\section{Tabelle und Verweise}
\label{sec:beispiel}

Die folgende Tabelle veranschaulicht den Grundaufbau.

\begin{table}[htbp]
  \centering
  \caption{Beispielhafte Kapitelstruktur}
  \label{tab:kapitel}
  \begin{tabular}{ll}
    \toprule
    Datei & Aufgabe \\
    \midrule
    \texttt{main.tex} & Hauptdokument \\
    \texttt{kapitel/einfuehrung.tex} & Erstes Kapitel \\
    \texttt{literatur.bib} & Literaturangaben \\
    \bottomrule
  \end{tabular}
\end{table}

Die Tabelle \autoref{tab:kapitel} kann mit ihrer Referenz angesprochen werden.

\clearpage
\section{Eine weitere Seite}

Diese zweite Kapitelseite dient als Test für das Scrollen in der mobilen
PDF-Vorschau. Auch umfangreichere Dokumente werden im Browser nur bei Bedarf
seitenweise gerendert.

Ein Bild lässt sich später mit \verb|\includegraphics| einfügen. Dabei sollte
die Grafik in einem Unterordner wie \texttt{bilder/} liegen.

\subsection{Änderungen veröffentlichen}

Nach einem erfolgreichen Build legt GitHub Actions die neue PDF gemeinsam
mit einer Versionsdatei auf GitHub Pages ab. Die geöffnete Vorschau überprüft
diese Datei regelmäßig und lädt das neue Dokument anschließend nach.
